\section*{Chapitre \ref{ch:s1:trend-following} --- Suivi de tendance}

\begin{solution}{exo:s1:trend-following:1}
$\lambda = 60/61 = 0.9836$~; les poids se réduisent de moitié après $\ln 0.5/\ln \lambda = 42$ jours.
\end{solution}

\begin{solution}{exo:s1:trend-following:2}
$0.4/(40 \times 0.25) = 0.04$, 4~\% du capital dans la matière première~; $0.4/(40 \times 0.06) = 0.167$, 16.7~\% dans le contrat à terme sur obligations. Chacun contribue pour 1~\% de
volatilité.
\end{solution}

\begin{solution}{exo:s1:trend-following:3}
Sur chaque cycle de hausse et de baisse, le mouvement total est nul, donc le premier terme, $\tfrac12(\sum r)^2$, est nul, et le portefeuille perd la moitié de la variance réalisée~: il
achète après la hausse et vend après la baisse. Les marchés heurtés représentent le coût du suiveur de tendance, comme l'érosion temporelle est celui du détenteur d'un straddle.
\end{solution}

\begin{solution}{exo:s1:trend-following:4}
$2.63 \times 0.09 + 0.03 \times (-0.30) - 0.01 = 0.218$~: environ 22~\% sur le trimestre, si la baisse est progressive, comme les krachs synthétiques.
\end{solution}

\begin{solution}{exo:s1:trend-following:5}
La règle à un mois fait tourner son encours brut 49 fois par an, et la \omterm{def:s1:trend-following:breakout}{règle de cassure} sur 250 jours 3 fois~; à quelques points de base par unité négociée, les coûts croissent avec
la rotation, tandis que l'edge brut de la règle rapide est faible. Le ratio de Sharpe de la cassure passe seulement de 0.80 à 0.78.
\end{solution}

\begin{solution}{exo:s1:trend-following:6}
Le ratio de Sharpe d'une décennie a une erreur type d'environ $1/\sqrt{10} = 0.32$, si bien que la plupart des décennies se situent entre $0.4 - 0.32 = 0.08$ et $0.4 + 0.32 = 0.72$. Une
décennie à 0.1 se situe dans cet intervalle~: prise seule, c'est un indice faible que l'edge a disparu. Chercher un mécanisme (corrélations, encombrement (crowding), coûts) avant
de conclure.
\end{solution}

\begin{solution}{exo:s1:trend-following:7}
La fenêtre de douze mois reste la meilleure (0.42), puis celle de six mois (0.39), puis celle de trois mois (0.13). Une fenêtre de trois mois suivrait la dérive la plus rapide, mais la mesurerait avec
beaucoup plus de bruit~: le rendement passé sur $L$ jours comporte une part de dérive croissant avec $L$ et une part de bruit croissant avec $\sqrt{L}$. À cette intensité de dérive, le
bruit domine les fenêtres courtes, si bien que la meilleure fenêtre reste longue même si les tendances sont plus rapides~; tous les ratios de Sharpe baissent.
\end{solution}

\begin{solution}{exo:s1:trend-following:8}
La meilleure parmi 200 essais corrélés est sélectionnée sur du bruit~: son ratio de Sharpe espéré est gonflé par le maximum de nombreuses estimations, chacune affectée d'une erreur
type d'environ 0.2 sur 25 ans. Il faut la déflater (livre~7, chapitre~20), ou mélanger les règles et juger le mélange~; hors échantillon (out-of-sample), il faut s'attendre à nettement moins.
\end{solution}

\begin{solution}{pb:s1:trend-following:1}
\begin{enumerate}
\item Négocier chaque marché dans le sens de son mouvement récent~; le signe du rendement passé sur $L$ jours~; une position aux nouveaux plus hauts ou plus bas sur $L$ jours, maintenue
jusqu'au signal contraire.
\item Rendement passé (poids égaux sur $L$ jours)~; croisement de moyennes mobiles (des poids en triangle, plus forts pour les jours récents)~; cassure (seulement les nouveaux extrêmes).
\item $\sum_t w_t r_t = \tfrac12(\sum_t r_t)^2 - \tfrac12\sum_t r_t^2$, en développant le carré.
\item Le carré des dérives persistantes et l'autocorrélation positive des rendements.
\item Dérives AR(1) de demi-vie d'un an et d'écart type égal à la moitié de la volatilité~; des carry dont la moitié est gagnée~; regroupement de volatilité par
classe d'actifs~; trois krachs de trois cents jours.
\item Signal multiplié par 40~\% sur 40 marchés, rapporté à une prévision de volatilité EWMA~; 2, 1, 2 et 4 points de base par unité négociée, selon la classe.
\item Rendement passé $-0.00$, 0.26, 0.66~; croisements 0.10, 0.43, 0.72~; cassures 0.17, 0.34, 0.78.
\item La dérive injectée est lente~; le mélange coûte l'écart avec la meilleure règle (0.55 pour neuf règles contre 0.78) et achète une assurance contre
le risque de choisir le mauvais rythme.
\item Dimensionnement sur une volatilité prévue~; 0.40 contre 0.31 à notionnel égal.
\item Il fait passer la volatilité de 7.7~\% à 10.2~\% et abaisse le ratio de Sharpe de 0.40 à 0.34~: il accroît son effet de levier lorsque les marchés divergent et que l'edge est faible.
\item Rendements positifs lors des baisses prolongées des actions~; $+2.3~\%$, $+21.1~\%$ et $+39.0~\%$ tandis que les actions perdaient 6.8~\%, 24.3~\% et 35.1~\%.
\item Les rendements trimestriels face aux actions ont une forme en U, de courbure 2.63~; le pire trimestre pour les actions a été le meilleur pour le portefeuille.
\item \textbf{Résultat nommé.} Lors des trois krachs synthétiques, le portefeuille de \omterm{def:s1:trend-following:trend}{suivi de tendance} à 10~\% de volatilité a rapporté $+2.3\%$, $+21.1\%$ et $+39.0\%$ contre des pertes sur les actions
de 6.8~\%, 24.3~\% et 35.1~\%~; le ratio de Sharpe après coûts croît avec la vitesse, de $-0.00$ (un mois) à 0.66 (douze mois) pour les rendements passés,
0.10 à 0.72 pour les croisements et de 0.17 à 0.78 pour les cassures.
\item 9.0~\% par an pour un ratio de Sharpe de 0.22, 1986--2024~; avec une erreur type d'environ 0.16, pas grand-chose.
\item Persistance d'un à douze mois sur 58 marchés et rendements anormaux substantiels~; rendements positifs dans chaque décennie depuis 1880 et bonne
performance dans huit crises sur dix.
\item Le niveau des corrélations deux à deux entre les marchés.
\item Une erreur type d'environ 0.32.
\item Un krach soudain, tel celui du 19 octobre 1987, plus rapide que la fenêtre d'observation des signaux.
\item Le chapitre~5 classe les actions les unes par rapport aux autres et est neutre vis-à-vis du marché~; ici chaque marché est négocié sur son propre passé, et le portefeuille peut être net acheteur ou net vendeur.
\item De grands mouvements, dans un sens comme dans l'autre, payés au prix de pertes sur les marchés heurtés.
\end{enumerate}
\end{solution}

\begin{solution}{iq:s1:trend-following:1}
Le \omterm{def:s1:trend-following:trend}{momentum temporel} (time-series momentum) négocie chaque actif sur son propre rendement passé et peut être net acheteur ou net vendeur sur l'ensemble du marché~; le \omterm{def:s1:momentum:cross}{momentum transversal} classe les actifs
les uns contre les autres, en achetant les gagnants et en vendant les perdants, et il est neutre par construction.
\end{solution}

\begin{solution}{iq:s1:trend-following:2}
Pour que chaque marché apporte un risque comparable~: sans cela, les marchés les plus volatils domineraient et le portefeuille ne serait pas diversifié~; et pour que le risque du portefeuille
soit plus stable à travers les régimes de volatilité.
\end{solution}

\begin{solution}{iq:s1:trend-following:3}
Comparer la perte à la distribution des rendements sur deux ans impliquée par le ratio de Sharpe espéré du portefeuille~; vérifier si le mécanisme a changé (corrélations,
encombrement, coûts, accès aux marchés)~; vérifier la mise en œuvre (roll, coûts, dimensionnement). Conserver la stratégie si la perte reste dans le bruit et que le mécanisme est intact.
\end{solution}

\begin{solution}{iq:s1:trend-following:4}
Roll des contrats à terme et série continue~; calendriers d'échéance~; marge au lieu de trésorerie~; multiplicateurs du contrat et pas de cotation~; rendements sur notionnel, et non sur
capital~; horaires de négociation et jours fériés propres à chaque marché.
\end{solution}

\begin{solution}{iq:s1:trend-following:5}
Il a couvert les baisses prolongées, car ses signaux deviennent vendeurs à mesure que la baisse se prolonge~; il ne couvre pas les krachs soudains, les retournements après une longue hausse, ni
les marchés heurtés où les actions et le portefeuille de \omterm{def:s1:trend-following:trend}{suivi de tendance} perdent tous deux.
\end{solution}

\begin{solution}{iq:s1:trend-following:6}
Si les rendements sont indépendants, $E[w_t r_t] = \sum_{k} E[r_{t-k}]E[r_t] = L\mu^2$. Pour un $\mu/\sigma$ petit, $\operatorname{Var}(w_t r_t) \approx E[w_t^2]E[r_t^2]
\approx L\sigma^2 \cdot \sigma^2$, donc le ratio de Sharpe quotidien vaut $L\mu^2/(\sqrt{L}\sigma^2) = \sqrt{L}(\mu/\sigma)^2$, à multiplier par $\sqrt{252}$ pour l'année~: il croît
avec la racine carrée de la fenêtre d'observation et le carré du ratio de Sharpe de la dérive.
\end{solution}
